% =====================================================================
%  3.6 Modellering — teori og oppgaver
%  Alle data i seksjonen er tenkte data.
% =====================================================================

\del{1}{Regresjon med digitale verktøy}

\begin{ordliste}
modell & En funksjon som beskriver en situasjon omtrentlig. & \mittsprak\\
regresjon & Å la et digitalt verktøy finne kurven av valgt type som ligger nærmest punktene. & \mittsprak\\
variabel & En størrelse som kan endre seg, for eksempel tid $x$ og antall $y$. & \mittsprak\\
modelltype & Lineær, eksponentiell, potens eller andregrad. & \mittsprak\\
\end{ordliste}

\begin{mustkunne}
\begin{minipage}[t]{0.48\linewidth}
\textbf{Fra data til modell}
\begin{enumerate}[leftmargin=1.3em,itemsep=0pt,topsep=2pt]
  \item Definer $x$ og $y$ med enhet.
  \item Plott punktene. Velg modelltype.
  \item Bruk regresjon i et digitalt verktøy.
  \item Tolk konstantene med enhet.
  \item Vurder: passer modellen med situasjonen? Er $R^2$ (forklaringsgraden) nær 1?
\end{enumerate}
Velg $x$ som «år etter 2010», ikke årstallet. Da er $x = 0$ første året, og konstanten $a$ blir verdien ved start.
\end{minipage}\hfill
\begin{minipage}[t]{0.49\linewidth}
\textbf{Hvilken modelltype?}\par\smallskip
\renewcommand{\arraystretch}{1.35}
\small
\begin{tabular}[t]{@{}l >{\raggedright\arraybackslash}p{4.1cm}@{}}
Lineær & like store tillegg i $y$ når $x$ øker med 1\\
Eksponentiell & $y$ ganges med samme tall når $x$ øker med 1\\
Potens & $y$ ganges med samme tall når $x$ dobles\\
Andregrad & har topp- eller bunnpunkt\\
\end{tabular}
\end{minipage}\medskip

Regresjon finner kurven som gjør avstandene til punktene totalt sett minst. Kurven treffer sjelden alle punktene.
\end{mustkunne}

\begin{eksempel}[1 — en eksponentiell modell]
Tenkte data: antall brukere (i tusen) av en ny app. $x$ er måneder etter lansering.

\smallskip
\begin{tabular}{@{}c|cccccc@{}}
$x$ & 0 & 1 & 2 & 3 & 4 & 5\\ \hline
$y$ & 2,1 & 3,0 & 4,3 & 6,2 & 9,0 & 12,7\\
\end{tabular}

\smallskip
\begin{stegtabell}
Differanser: $0{,}9;\ 1{,}3;\ 1{,}9;\ 2{,}8;\ 3{,}7$ & Differansene vokser. Da er det ikke en lineær modell.\\
Forhold: $\frac{3{,}0}{2{,}1} \approx 1{,}43$, $\frac{4{,}3}{3{,}0} \approx 1{,}43$, $\frac{6{,}2}{4{,}3} \approx 1{,}44$ & Nesten samme faktor hver gang. Velg eksponentiell: $f(x) = a \cdot b^x$.\\
$f(x) \approx 2{,}09 \cdot 1{,}44^x$ & Bruk eksponentiell regresjon i et digitalt verktøy.\\
$a \approx 2{,}09$ og $b \approx 1{,}44$ & $a$: omtrent 2\,090 brukere ved lansering. $b$: antallet ganges med $1{,}44$ per måned, altså $44\,\%$ vekst per måned.\\
Vurdering: $f(5) \approx 12{,}8$ & Målt verdi er $12{,}7$. Modellen passer, og $R^2$ er nær 1.\\
\end{stegtabell}
\end{eksempel}

\begin{eksempel}[2 — velg $x$ med omhu]
Tenkte data: medlemmer i et turlag ved årsslutt 2010 til 2015: 212, 228, 247, 261, 280, 297. Tilleggene (16, 19, 14, 19, 17) er nesten like, så vi velger lineær modell.

\begin{stegtabell}
$x = $ årstall & Regresjon gir konstantledd $\approx -33\,958$. Meningsløst: «$-33\,958$ medlemmer i år 0».\\
$x = $ år etter 2010 & Regresjon gir $f(x) \approx 17{,}0x + 212$. Nå er $x = 0$ i 2010.\\
Tolking & $17{,}0$: 17 nye medlemmer per år. $212$: medlemmer i 2010.\\
\end{stegtabell}
\end{eksempel}

\begin{vanligfeil}
\begin{itemize}[leftmargin=1.2em,itemsep=1pt,topsep=0pt]
  \item «Bruk årstallet som $x$» gir store, meningsløse konstanter, som $-33\,958$ i eksempel 2. Bruk år etter startåret.
  \item «$b = 1{,}44$ betyr 1,44 brukere» er feil. Riktig: antallet ganges med $1{,}44$ per måned. Skriv alltid enhet.
  \item «Alt som øker, er eksponentielt» gjelder ikke generelt. Tabellen $4, 7, 10, 13$ øker, men med like store tillegg ($+3$). Den er lineær.
\end{itemize}
\end{vanligfeil}

\nivaahode{1}

\oppg{1.1}{1}{\uh}
Tenkte data. Hvilken modelltype passer: lineær, eksponentiell eller potens? Begrunn med tallene.
\begin{center}
\textbf{a)}\ \begin{tabular}[t]{c|cccc} $x$ & 1 & 2 & 3 & 4\\ \hline $y$ & 5 & 9 & 13 & 17\end{tabular}\qquad
\textbf{b)}\ \begin{tabular}[t]{c|cccc} $x$ & 0 & 1 & 2 & 3\\ \hline $y$ & 5 & 15 & 45 & 135\end{tabular}\qquad
\textbf{c)}\ \begin{tabular}[t]{c|cccc} $x$ & 1 & 2 & 4 & 8\\ \hline $y$ & 5 & 20 & 80 & 320\end{tabular}
\end{center}

\oppg{1.2}{1}{\uh}
Tenkte modeller. Tolk tallene med enhet.
\begin{tasks}(1)
  \task Et basseng fylles: $V(t) = 45t + 120$. Her er $V$ liter og $t$ minutter. Hva forteller $45$ og $120$?
  \task En konto har $N(x) = 800 \cdot 1{,}15^x$ følgere. Her er $x$ uker. Hva forteller $800$ og $1{,}15$?
  \task Regn ut $V(10)$ og $N(2)$.
\end{tasks}

\oppg{1.3}{1}{\mh}
Tenkte data: antall elbiler i en kommune. $x$ er år etter 2015.

\smallskip
\begin{tabular}{@{}c|cccccc@{}}
$x$ & 0 & 1 & 2 & 3 & 4 & 5\\ \hline
$y$ & 120 & 171 & 226 & 272 & 331 & 379\\
\end{tabular}

\smallskip
Bruk et digitalt verktøy. Fullfør.
\begin{tasks}(1)
  \task Plott punktene. Ser de ut til å ligge langs en rett linje?
  \task Bruk lineær regresjon: $f(x) = $ \skrivlinje[2cm]$\,x\ +$ \skrivlinje[2cm]
  \task Antall elbiler øker med omtrent \skrivlinje[2cm] per \skrivlinje[2cm]
\end{tasks}

\nivaahode{2}

\oppg{1.4}{2}{\mh}
Tenkte data: antall visninger av en video per dag. $x$ er dager etter publisering.

\smallskip
\begin{tabular}{@{}c|ccccccc@{}}
$x$ & 0 & 1 & 2 & 3 & 4 & 5 & 6\\ \hline
$y$ & 120 & 160 & 215 & 290 & 380 & 510 & 690\\
\end{tabular}

\begin{tasks}(1)
  \task Regn ut differanser og forhold mellom nabopunktene. Hvilken modelltype passer?
  \task Bruk regresjon og skriv modellen.
  \task Tolk $a$ og $b$ med enhet. Hvor mange prosent øker antallet per dag?
  \task Hvor mange visninger er det dag 8 ifølge modellen? Svar med to gjeldende siffer.
\end{tasks}

\oppg{1.5}{2}{\mh}
Tenkte data: svingetiden $T$ (sekunder) til et pendel med lengde $L$ (meter).

\smallskip
\begin{tabular}{@{}c|ccccc@{}}
$L$ & 0,25 & 0,50 & 1,00 & 1,50 & 2,00\\ \hline
$T$ & 1,00 & 1,42 & 2,01 & 2,46 & 2,84\\
\end{tabular}

\begin{tasks}(1)
  \task Regn ut $\dfrac{T(0{,}50)}{T(0{,}25)}$ og $\dfrac{T(1{,}00)}{T(0{,}50)}$. Hva forteller det om modelltypen?
  \task Bruk regresjon og skriv modellen.
  \task Finn svingetiden når $L = 3{,}0$ m.
  \task Hvor lang må pendelen være for at svingetiden skal bli $3{,}0$ s? Svar med to gjeldende siffer.
\end{tasks}

\oppg{1.6}{2}{\uh}
Tenkte data: $x = 0, 1, 2, 3, 4$ og $y = 50, 60, 72, 86, 104$. En elev vil finne modellen og regne ut $y$ for $x = 5$.
Finn \textbf{den første feilen}. Regn riktig.

\smallskip
\begin{tabular}{@{}l p{14cm}@{}}
Steg 1 & Differansene er $10, 12, 14, 18$. De er ikke like, så modellen er ikke lineær.\\
Steg 2 & Forholdene er $1{,}2;\ 1{,}2;\ 1{,}19;\ 1{,}21$. De er omtrent like, så modellen er eksponentiell.\\
Steg 3 & Regresjon gir $f(x) \approx 50 \cdot 1{,}2^x$.\\
Steg 4 & $f(5) = 50 \cdot 1{,}2 \cdot 5 = 300$.\\
\end{tabular}

\nivaahode{3}

\oppg{1.7}{3}{\uh}
Er påstanden alltid, noen ganger eller aldri sann? Begrunn svaret.
\begin{tasks}(1)
  \task Når $y$ øker med like mye hver gang $x$ øker med 1, passer en lineær modell nøyaktig.
  \task Kurven fra regresjon går gjennom alle datapunktene.
  \task Stigningstallet i en lineær modell er det samme om vi måler $x$ i år eller i måneder.
\end{tasks}

\oppg{1.8}{3}{\mh}
Lag selv en tabell med fem punkter, $x = 0, 1, 2, 3, 4$, der en eksponentiell modell passer mye bedre enn en lineær.
Bruk regresjon og sammenlign $R^2$ for de to modellene. Forklar hvordan du valgte tallene.

\oppg{1.9}{3}{\uh}
Tenkte data: $x = 1, 2, 4, 8$ og $y = 6, 12, 24, 48$. Ola ser at $y$ dobles for hvert hopp og sier at modellen er eksponentiell.
\begin{tasks}(1)
  \task Forklar hvorfor Ola tar feil.
  \task Finn en riktig modell.
\end{tasks}

\begin{sjekk}
\begin{tasks}(1)
  \task Tabellen har $x = 0, 1, 2, 3$ og $y = 80, 40, 20, 10$. Hvilken modelltype passer?
  \task Hvorfor bruker vi «år etter 2010» som $x$ og ikke årstallet?
  \task $B(t) = 3{,}5t + 12$ er batteriprosent etter $t$ minutter lading. Hva forteller $3{,}5$ og $12$?
\end{tasks}
\end{sjekk}

% ---------------------------------------------------------------------

\del{2}{Vurdere modeller og lese tekster}

\begin{ordliste}
gyldighetsområde & De verdiene av $x$ der modellen gir mening. & \mittsprak\\
prognose & Et anslag om fremtiden regnet ut med en modell. & \mittsprak\\
forklaringsgrad & Tallet $R^2$ mellom 0 og 1. Nær 1 betyr at kurven ligger tett inntil punktene. & \mittsprak\\
vurdere & Å sjekke om modellen og svarene gir mening i situasjonen. & \mittsprak\\
\end{ordliste}

\begin{mustkunne}
\begin{minipage}[t]{0.50\linewidth}
\textbf{Hvor langt kan vi stole på modellen?}
\begin{itemize}[leftmargin=1.2em,itemsep=1pt,topsep=2pt]
  \item Innenfor dataene er modellen brukbar. Utenfor er en prognose usikker, og mer usikker jo lenger ut vi går.
  \item $R^2$ nær 1 sier at kurven passer med punktene. Det sier ikke at modellen passer med situasjonen.
  \item Sjekk grensene: negative verdier? Over $100\,\%$? Vekst for alltid?
  \item Skriv «ifølge modellen», ikke «modellen viser at».
\end{itemize}
\end{minipage}\hfill
\begin{minipage}[t]{0.46\linewidth}
\vspace{0pt}\centering
\begin{tikzpicture}[baseline=(current bounding box.north)]
\begin{axis}[width=\linewidth, height=0.72\linewidth, axis lines=middle, xmin=0, xmax=24.5, ymin=-14, ymax=20,
  xtick={6,12,20}, ytick={-10,0,10}, tick label style={font=\scriptsize}, xlabel={$x$}, ylabel={$y$}, label style={font=\small}]
\addplot[domain=6:20, thick, hoved, samples=40] {-0.15685*x^2 + 4.3006*x - 12.682};
\addplot[domain=0:6, thick, red!70!black, dashed, samples=20] {-0.15685*x^2 + 4.3006*x - 12.682};
\addplot[domain=20:24, thick, red!70!black, dashed, samples=20] {-0.15685*x^2 + 4.3006*x - 12.682};
\addplot[only marks, mark=*, mark size=1.5pt] coordinates {(6,7.8) (8,11.5) (10,14.1) (12,16.4) (14,17.0) (16,16.3) (18,13.9) (20,10.4)};
\end{axis}
\end{tikzpicture}\par
{\footnotesize Heltrukket: der dataene er. Stiplet: utenfor dataene.}
\end{minipage}
\end{mustkunne}

\begin{eksempel}[1 — hvor gjelder modellen?]
Tenkte data: temperatur ($^\circ$C) målt hver andre time en sommerdag. $x$ er klokketimer.

\smallskip
\begin{tabular}{@{}c|cccccccc@{}}
$x$ & 6 & 8 & 10 & 12 & 14 & 16 & 18 & 20\\ \hline
$y$ & 7,8 & 11,5 & 14,1 & 16,4 & 17,0 & 16,3 & 13,9 & 10,4\\
\end{tabular}

\smallskip
\begin{stegtabell}
$T(x) \approx -0{,}157x^2 + 4{,}30x - 12{,}7$ & Andregradsregresjon: temperaturen stiger og synker, så vi trenger et toppunkt. $R^2 \approx 0{,}99$.\\
$T(13) \approx 16{,}7$ & Innenfor dataene ($6 \le x \le 20$). Brukbar.\\
$T(2) \approx -4{,}7$ & Utenfor dataene. En tenkt måling kl. 02 viste $8{,}5\,^\circ$C. Modellen bommer med $13{,}2$ grader.\\
$T(0) \approx -12{,}7$ & Parabelen synker for alltid til venstre. Slik er ikke en natt.\\
Konklusjon & Gyldighetsområde: $6 \le x \le 20$. Ifølge modellen er det varmest omtrent kl. 14.\\
\end{stegtabell}
\end{eksempel}

\begin{eksempel}[2 — lese en tekst]
Tenkt avisutdrag: «Svømmehallen i Bakkeby hadde 12\,000 besøkende i 2020 og 15\,600 i 2023.
– Hvis veksten fortsetter, har vi over 100\,000 besøkende i 2040, sier daglig leder.»

\begin{stegtabell}
Tall og påstand & Tall: $12\,000$ (2020) og $15\,600$ (2023). Påstand: over $100\,000$ i 2040.\\
$\dfrac{15\,600}{12\,000} = 1{,}3$ & Økning på $30\,\%$ på tre år.\\
Lineær: $\dfrac{3\,600}{3} = 1\,200$ & Samme tillegg per år. Da blir 2040: $12\,000 + 1\,200 \cdot 20 = 36\,000$.\\
Eksponentiell: $12\,000 \cdot 1{,}3^{\tfrac{20}{3}} \approx 69\,000$ & Faktor $1{,}3$ per tre år. 20 år er $\frac{20}{3}$ perioder.\\
Vurdering & Ingen av modellene gir 100\,000. To tall kan ikke avgjøre modelltypen. Tallene i teksten støtter ikke påstanden.\\
\end{stegtabell}
\end{eksempel}

\begin{vanligfeil}
\begin{itemize}[leftmargin=1.2em,itemsep=1pt,topsep=0pt]
  \item «Modellen viser at det blir 69\,000 besøkende» er for sterkt. Skriv «ifølge modellen» og si hva som må være sant.
  \item «$R^2 = 1$, så modellen er riktig» gjelder ikke generelt. En rett linje gjennom to punkter har alltid $R^2 = 1$, men to punkter sier lite.
  \item «$N(x) = -3x + 45$ gir $N(20) = -15$ brukere» er en modell brukt utenfor gyldighetsområdet. Antall kan ikke være negativt.
\end{itemize}
\end{vanligfeil}

\nivaahode{1}

\oppg{2.1}{1}{\uh}
Tenkt modell: En vare har tilbud. Prisen er $P(x) = -2x + 60$ kroner, der $x$ er dager etter at tilbudet startet.
Modellen er laget ut fra prisene de ti første dagene.
\begin{tasks}(1)
  \task Regn ut $P(4)$.
  \task Regn ut $P(35)$. Hva er galt med svaret?
  \task Foreslå et gyldighetsområde for modellen.
\end{tasks}

\oppg{2.2}{1}{\uh}
Skriv om hver setning slik at den ikke sier mer enn modellen kan støtte.
\begin{tasks}(1)
  \task «Modellen viser at folketallet blir 9\,000 i 2040.»
  \task «Regresjonen beviser at prisen dobles hvert femte år.»
  \task «Siden $R^2 = 0{,}99$ er modellen riktig.»
\end{tasks}

\oppg{2.3}{1}{\uh}
Svar kort.
\begin{tasks}(1)
  \task Hva sier $R^2 = 1$ om punktene og kurven?
  \task Hva sier $R^2 = 0{,}12$ om modellen?
  \task En modell har $R^2 = 0{,}6$. Hva bør du gjøre før du bruker den?
\end{tasks}

\nivaahode{2}

\oppg{2.4}{2}{\mh}
Tenkte data: antall besøkende per time i en butikk som åpner kl. 10. $x$ er timer etter kl. 10.

\smallskip
\begin{tabular}{@{}c|cccccccc@{}}
$x$ & 0 & 1 & 2 & 3 & 4 & 5 & 6 & 7\\ \hline
$y$ & 12 & 25 & 33 & 38 & 36 & 31 & 22 & 9\\
\end{tabular}

\begin{tasks}(1)
  \task Bruk andregradsregresjon og skriv modellen.
  \task Når er det flest besøkende ifølge modellen?
  \task Regn ut $f(8)$. Hva er galt med svaret?
  \task Angi et gyldighetsområde og begrunn det.
\end{tasks}

\oppg{2.5}{2}{\mh}
Tenkt avisutdrag: «Fjellvik vokser. Kommunen hadde 5\,200 innbyggere i 2015 og 5\,640 i 2019. Ved årsskiftet 2023 var tallet 6\,080.
– Veksten er over 20 prosent siden 2015, og i 2040 har vi 8\,500 innbyggere, sier ordføreren.»
\begin{tasks}(1)
  \task Sett tallene i en tabell. La $x$ være år etter 2015.
  \task Regn ut den prosentvise veksten fra 2015 til 2023. Stemmer påstanden «over 20 prosent»?
  \task Lag en lineær modell og finn folketallet i 2040.
  \task Hvor mange flere innbyggere per år enn i modellen trengs for å nå 8\,500 i 2040? Hva må ordføreren anta?
\end{tasks}

\oppg{2.6}{2}{\uh}
Tenkt modell: Månedsprisen for en tjeneste er $P(x) = 2{,}5x + 40$ kroner, der $x$ er år etter 2020.
En elev vil finne når prisen er 100 kroner. Finn \textbf{den første feilen}. Regn riktig.

\smallskip
\begin{tabular}{@{}l p{14cm}@{}}
Steg 1 & $2{,}5x + 40 = 100$\\
Steg 2 & $2{,}5x = 60$\\
Steg 3 & $x = \dfrac{60}{2{,}5} = 25$\\
Steg 4 & $2020 + 25 = 2045$\\
Steg 5 & «Modellen viser at prisen blir 100 kroner i 2045.»\\
\end{tabular}

\nivaahode{3}

\oppg{2.7}{3}{\uh}
Er påstanden alltid, noen ganger eller aldri sann? Begrunn svaret.
\begin{tasks}(1)
  \task Modellen med høyest $R^2$ gir den beste prognosen.
  \task En lineær modell med negativt stigningstall gir til slutt negative $y$-verdier.
  \task Bytter vi fra årstall til år etter 2010 som $x$ i en eksponentiell modell, endres vekstfaktoren $b$.
\end{tasks}

\oppg{2.8}{3}{\mh}
Tenkte data: antall kunder (i tusen) hos en fiberleverandør. $x$ er år etter oppstart.

\smallskip
\begin{tabular}{@{}c|cccccc@{}}
$x$ & 0 & 1 & 2 & 3 & 4 & 5\\ \hline
$y$ & 100 & 111 & 124 & 137 & 152 & 168\\
\end{tabular}

\begin{tasks}(1)
  \task Lag en lineær og en eksponentiell modell med regresjon. Finn $R^2$ for begge.
  \task Bruk begge modellene til å anslå antall kunder ved $x = 15$.
  \task Begge modellene passer godt. Hvilken bør leverandøren bruke til en prognose over ti år? Argumenter ut fra situasjonen.
\end{tasks}

\oppg{2.9}{3}{\mh}
Tenkt eksempel: Et land har 5 millioner innbyggere. En eksponentiell modell sier at folketallet vokser med $1{,}5\,\%$ per år.
\begin{tasks}(1)
  \task Hva sier modellen om folketallet om 200 år?
  \task Forklar hvorfor modellen ikke kan gjelde så lenge.
  \task Beskriv hvordan en kurve som passer bedre, ville sett ut.
\end{tasks}

\begin{sjekk}
\begin{tasks}(1)
  \task Modellen $N(x) = -4x + 50$ gir antall billetter $N$ etter $x$ dager, med $x \ge 0$. For hvilke $x$ er $N(x) \ge 0$?
  \task Skriv om: «Modellen viser at det blir 500 brukere i 2030.»
  \task Hvorfor sier høy $R^2$ lite om en prognose langt utenfor dataene?
\end{tasks}
\end{sjekk}
